\section{Síntesis y ampliación}

\subsection{Conclusiones centrales}

\begin{enumerate}
\item Coding es el segundo acto de la AGI, porque lleva al modelo de conversar a ejecutar tareas en el mundo digital.
\item La ventana de competencia entre las empresas de modelos de frontera se está reordenando; el éxito de la etapa anterior puede convertirse en dependencia de la trayectoria en la siguiente.
\item La Harness Engineering determina si el modelo puede entrar en flujos de trabajo reales.
\item El modelo se vuelve la pieza central de una nueva generación de sistemas operativos porque puede coordinar herramientas, contexto, permisos y estado de las tareas.
\item Coding/Agent traerá una deflación del trabajo de oficina y una ventana de desempleo; la sociedad podría absorber el cambio más despacio de lo que se difunde la tecnología.
\end{enumerate}

\subsection{Preguntas abiertas}

\begin{itemize}
\item ¿Llegarán los modelos de Coding a ser, como la GPU, el acelerador básico de toda la investigación en IA?
\item ¿El sistema operativo basado en modelos quedará en manos de las grandes empresas actuales o de una nueva generación de empresas de aplicaciones AI-native?
\item ¿En qué industrias se producirá primero la deflación del trabajo de oficina y cómo redistribuirá la sociedad los beneficios de la productividad?
\item ¿Podrán las empresas chinas de modelos alcanzar el segundo acto apoyándose en sus costos, sus aplicaciones y su ecosistema de desarrolladores?
\end{itemize}

\subsection{Lecturas adicionales}

\begin{itemize}
\item EP139, historia técnica de los Agent: para entender por qué Coding/Agent forma parte de una genealogía tecnológica más amplia.
\item EP138, entrevista con Luo Fuli (罗福莉): para entender cómo el paradigma Agent cambia el posentrenamiento y la asignación de capacidad de cómputo.
\item Práctica con herramientas como Claude Code, Codex y OpenClaw: para observar cómo los modelos de Coding cambian el flujo de trabajo de investigación y desarrollo.
\item Harness engineering, agent evaluation, computer use agent benchmark: para entender la capa intermedia por la que los modelos entran en los sistemas de producción.
\end{itemize}

\begin{importantbox}{Juicio final}
Si en 2023 la pregunta era «¿el modelo puede hablar?», en 2026 la pregunta es «¿el modelo puede trabajar?». Coding es el primer terreno firme de esta transición; quien logre conectarlo a flujos de trabajo reales estará más cerca de la próxima generación de sistemas operativos de IA.
\end{importantbox}

\end{document}
